\documentclass[aspectratio=169,11pt]{beamer}
\usepackage[utf8]{inputenc}
\usepackage[T1]{fontenc}
\usepackage[spanish]{babel}
\usepackage{lmodern}
\usepackage{booktabs}
\usepackage{tabularx}
\usepackage{hyperref}
\usepackage{enumitem}

\usetheme{Madrid}
\usecolortheme{default}
\setbeamertemplate{navigation symbols}{}
\setbeamertemplate{footline}[frame number]

\title{IA agéntica para la creación de contenido académico}
\subtitle{Sesión 4 --- Libros/ebooks II: secciones, ejercicios y consistencia}
\author{Juliho Castillo}
\institute{CADi 40\,h · Escuela de Ingeniería y Ciencias\\
Tecnológico de Monterrey}
\date{Clave sugerida: \texttt{CADI-EIC-IA-AGENTICA-01}}

\begin{document}

\begin{frame}
\titlepage
\end{frame}

\begin{frame}{Agenda (3\,h) y producto E4}
\begin{tabularx}{\textwidth}{@{}lclX@{}}
\toprule
Bloque & Tiempo & Contenido \\
\midrule
Recap outlines & 10 min & Priorizar secciones \\
Demo sección & 25 min & Restricciones y regeneración \\
Figuras y ejercicios & 20 min & Alineación a objetivos \\
Taller secciones & 55 min & Generación + edición \\
Consistencia & 35 min & Agente de estilo/glosario \\
Cierre & 15 min & Diff v0$\rightarrow$v1; TI \\
\bottomrule
\end{tabularx}
\vspace{0.5em}
\textbf{E4:} 1--2 secciones + artefacto didáctico + reporte de consistencia.
\end{frame}

\begin{frame}{Resultados de aprendizaje (sesión)}
Al finalizar, usted podrá:
\begin{enumerate}
\item Generar de forma supervisada \textbf{1--2 secciones} coherentes en tono y nivel
\item Producir al menos un \textbf{ejercicio o figura} alineado a objetivos
\item Ejecutar un \textbf{pase de consistencia} (glosario, notación, unidades, estilo)
\item Documentar un \textbf{diff narrado} v0$\rightarrow$v1 (qué cambió el humano)
\end{enumerate}
\end{frame}

\begin{frame}{De outline a sección corta}
\begin{itemize}
\item Elija 1--2 nodos prioritarios del E3
\item Genere \textbf{por secciones cortas} (reduce alucinaciones)
\item Edite humanamente; marque \texttt{[VERIFICAR]}
\item No intente el capítulo completo en un solo prompt
\end{itemize}
\begin{block}{CEDDIE}
No improvisar normativa de integridad / uso de IA --- verificar por edición.
\end{block}
\end{frame}

\begin{frame}{Demo --- restricciones y el ``no'' humano}
Modelo en vivo:
\begin{enumerate}
\item Prompt con brief + nodo + notación + nivel
\item Primera salida: detectar relleno o fórmula dudosa
\item Humano dice \textbf{no}: regenerar o editar
\item Conservar v0 y anotar decisión
\end{enumerate}
\begin{alertblock}{Mensaje clave}
Generar por secciones cortas $>$ un muro de texto.
\end{alertblock}
\end{frame}

\begin{frame}{Figuras, tablas y ejercicios STEM}
\begin{itemize}
\item Nacen de \textbf{objetivos / RA}, no de relleno
\item Cada ejercicio: enunciado + solución o rúbrica breve
\item Figuras: propósito didáctico + leyenda + fuentes si aplica
\item Tablas: unidades visibles; sin datos inventados
\end{itemize}
Si no hay fuente: \texttt{[VERIFICAR]} o no incluir.
\end{frame}

\begin{frame}{Agente de consistencia: qué sí / qué no}
\begin{tabularx}{\textwidth}{@{}XX@{}}
\toprule
\textbf{Sí (útil)} & \textbf{No (límites)} \\
\midrule
Términos vs.\ glosario & Sustituir revisión disciplinar \\
Notación / unidades & Validar teoremas nuevos \\
Estilo y tono & Fabricar citas ``faltantes'' \\
Duplicados / contradicciones & Decidir pedagogía del curso \\
\bottomrule
\end{tabularx}
\vspace{0.8em}
El agente amplifica; el humano certifica.
\end{frame}

\begin{frame}{Plantilla de reporte de consistencia}
Cinco a ocho líneas:
\begin{itemize}
\item Hallazgos (alto / medio / bajo)
\item Correcciones aplicadas
\item Pendientes \texttt{[VERIFICAR]}
\item Cambios de notación o glosario
\item Qué rechazó el humano
\end{itemize}
\end{frame}

\begin{frame}{Diff narrado v0$\rightarrow$v1}
Documente:
\begin{enumerate}
\item Qué generó la IA (v0)
\item Qué cambió / cortó / reordenó la persona
\item Por qué (rigor, nivel, tono, integridad)
\item Qué queda abierto
\end{enumerate}
Hábito que se cierra en el agent log (sesiones 7--8).
\end{frame}

\begin{frame}{Integridad STEM en secciones}
Checklist rápido:
\begin{itemize}
\item[$\square$] Ecuaciones y símbolos coherentes con el mapa
\item[$\square$] Unidades sin mezclar sistemas
\item[$\square$] Ninguna cita inventada
\item[$\square$] Afirmaciones empíricas con fuente o \texttt{[VERIFICAR]}
\end{itemize}
\end{frame}

\begin{frame}{Actividad de taller}
\textbf{Entregable E4:}
\begin{enumerate}
\item \textbf{Parte A --- Secciones} (30 min): 1--2 nodos; editar; marcar \texttt{[VERIFICAR]}
\item \textbf{Parte B --- Ejercicio/figura} (15 min): alineado a un RA
\item \textbf{Parte C --- Consistencia} (15 min): pase + 5--8 líneas de hallazgos
\end{enumerate}
Nombre: \texttt{Apellido\_Sesion4\_secciones-consistencia}
\end{frame}

\begin{frame}{Rúbrica rápida (E4)}
\begin{tabularx}{\textwidth}{@{}lXX@{}}
\toprule
Criterio & Bien & A mejorar \\
\midrule
Alineación EIC & Nivel y notación claros & Genérico \\
HITL & Diff / \texttt{[VERIFICAR]} visibles & Aceptación pasiva \\
Artefacto & Ejercicio/figura con RA & Relleno decorativo \\
Consistencia & Reporte concreto & ``Se ve bien'' \\
\bottomrule
\end{tabularx}
\end{frame}

\begin{frame}{Anti-patrones}
\begin{itemize}
\item Generar el capítulo entero de un golpe
\item Figuras sin propósito didáctico
\item Agente de consistencia como ``aprobación automática''
\item Borrar v0 sin documentar el diff
\item Inventar bibliografía para ``rellenar''
\end{itemize}
\end{frame}


\begin{frame}{Workshop cue --- calidad $>$ volumen}
\begin{itemize}
\item Quien termine la primera sección: segundo ejercicio + tabla
\item Circulación: notación, tono, \texttt{[VERIFICAR]}
\item Si la herramienta falla: continuar en documento colaborativo
\end{itemize}
\end{frame}

\begin{frame}{Puente al bloque de presentaciones}
\begin{itemize}
\item Hoy cerramos prosa de libro/ebook
\item Sesiones 5--6: narrativa, deck y accesibilidad
\item Su sección E4 puede alimentar el storyboard
\end{itemize}
\end{frame}

\begin{frame}{Takeaways de la sesión}
\begin{enumerate}
\item Secciones cortas + edición humana
\item Artefactos nacen de RA, no de relleno
\item Consistencia $\neq$ revisión disciplinar completa
\item Diff narrado = transparencia del proceso
\end{enumerate}
\end{frame}

\begin{frame}{Trabajo independiente}
Antes de la sesión 5:
\begin{itemize}
\item Pulir sección (bloque C)
\item Practicar plantillas (bloque B)
\item Traer idea de \textbf{unidad o conferencia} a narrar
\end{itemize}
Parte de las \textbf{16\,h} independientes del diseño CADi.
\end{frame}

\begin{frame}{Prep sesión 5 --- Presentaciones I}
Próxima sesión (3\,h):
\begin{itemize}
\item De objetivos a mensajes clave y arco narrativo
\item Storyboard 8--12 nodos
\item Anticipar malentendidos STEM
\end{itemize}
Traiga E4 (o la sección pulida) como materia prima.
\end{frame}

\begin{frame}{Gracias}
\begin{center}
{\Large Juliho Castillo}\\[0.5em]
Escuela de Ingeniería y Ciencias\\[1em]
CADi · IA agéntica para la creación de contenido académico\\[1.5em]
{\small \emph{Verificar políticas de integridad y uso de IA con CEDDIE antes de cada edición.}}
\end{center}
\end{frame}

\end{document}
